% ==========================================
% 第四篇 导入过场页
% ==========================================
\cleardoublepage
\thispagestyle{empty}

\vspace*{6cm}

\begin{center}
	{\heiti\erhao 第四篇\quad 我们是谁} 
	\addcontentsline{toc}{part}{第四篇\quad 我们是谁}
	
	\vspace{2cm} 
	
	{\songti\xiaosi \textit{“一支理解密码学底层逻辑、懂得安防行业落地规则、\\ [0.5em] 并在量子安全时代窗口期果断入局的顶尖产学研团队。”}}
\end{center}

\vfill
\cleardoublepage

% ==========================================
% 第十章
% ==========================================
\chapter{团队与融资计划}

本项目的核心壁垒不仅在于前沿的技术架构，更源于我们拥有国内顶尖的密码学产学研落地团队。依托武汉大学国家网络安全学院及计算机学院的深厚学术底蕴，我们构建了一支由国家杰青、重点研发首席科学家领衔，由具备极客精神和商业敏锐度的优秀本科生主导的复合型战队。

\section{核心创始团队}

学生创始团队来自武汉大学密码学与计算机顶尖拔尖班级，兼具底层代码硬实力与商业架构规划能力。

\begin{itemize}
	\item \textbf{张睿恒 | 团队负责人 / 首席执行官 (CEO)} \\
	国家网络安全学院密码试验班 2024级本科生。负责项目的整体战略统筹、底层后量子算法库代码开发与门禁协议适配。曾斩获第十九届“挑战杯”全国大学生课外学术科技作品竞赛\textbf{国家级三等奖}。
	
	\item \textbf{傅承阳 | 核心技术骨干 / 首席技术官 (CTO)} \\
	国家网络安全学院密码试验班 2024级本科生。主导PQC算法的工程化实现与性能极限优化。具备极强的代码落地能力，荣获 \textbf{2025年雷军计算机奖学金}，及2025年“挑战杯”\textbf{国家级三等奖}。
	
	\item \textbf{胡斯舜 | 战略与商业负责人 (COO)} \\
	国家网络安全学院密码科学与技术试验班 2024级本科生。主导项目商业模式设计、财务预测构建及商业计划书撰写，精准把控抗量子基建市场的商业化演进路径。
	
	\item \textbf{汪扬帆 | 商务与路演负责人 (CMO)} \\
	国家网络安全学院密码试验班 2025级本科生。专业绩点稳居\textbf{年级前10\%}，负责对外路演答辩与客户痛点需求挖掘，具备扎实的密码学理论基础与卓越的表达感染力。
	
	\item \textbf{李恺睿 | 运营与合规负责人} \\
	国家网络安全学院密码试验班 2024级本科生。负责项目申报书矩阵构建、知识产权规划与国家密码政策合规性梳理。
	
	\item \textbf{刘芸萱、王一茗 | 品牌视觉与市场宣传总监} \\
	分别来自国家网络安全学院计算机类2025级与新闻与传播学院新闻系2024级（曾获甲等奖学金）。负责公司品牌VI设计、产品UI交互呈现及高质量发布会级路演视觉制作，赋予硬核科技产品卓越的商业美感。
\end{itemize}

\section{顶级科学家顾问委员会}

我们的指导教师团队由国内密码学领域的领军人物组成，为本项目的底层PQC算法、国密合规以及行业标准制定提供了\textbf{无法复制的降维打击能力与学术护城河}。

\begin{description}[labelindent=0pt, leftmargin=1em, style=sameline]
	\item[\textbf{何德彪 教授（首席学术顾问）：}] 武汉大学计算机学院教授、院长，\textbf{国家杰出青年科学基金获得者}。主持国家自然科学基金、科技部重点研发计划等国家级科研项目。近5年，在 IEEE TIFS, IEEE TDSC, ASIACRYPT, PKC 等国际密码学顶级期刊/会议上发表论文40余篇。
	
	\item[\textbf{彭聪 副教授（PQC标准推进顾问）：}] 武汉大学国家网络安全学院副教授，\textbf{国家重点研发计划青年项目首席科学家}。在 ASIACRYPT 等著名期刊发表论文70余篇。其参与实现的“SM2MLKEM768”近期成功获得 \textbf{IANA TLS官方编号：4590}，这是\textbf{我国在后量子密码领域首次获得该项国际认可}，直接赋能本项目底层算法的国际领先性。
	
	\item[\textbf{冯琦 副教授（密码行业标准顾问）：}] 武汉大学国家网络安全学院副教授。发表顶级学术论文40余篇，申请发明专利60余项。\textbf{参与制定/修订密码行业标准2项、成果入选密码行业标准草案}，为项目未来的合规落地与国标主导权奠定坚实基础。
	
	\item[\textbf{李小红 副教授（系统工程顾问）：}] 武汉大学计算机学院软件工程系副教授，曾在美国 Purdue 大学访问。先后主持8个横向课题及多项国家级/国防保密课题，为门禁系统工程化落地与军工级保密需求提供指导。
	
	\item[\textbf{罗敏 教授（网络安全生态顾问）：}] 武汉大学国家网络安全学院副教授、博士生导师。密码科学与技术系党支部书记，CCF区块链专业委员会委员。为项目的物联网生态拓展与分布式安全架构提供前瞻性指引。
\end{description}

\section{融资需求与资金用途}

基于团队现状与规模化扩张需求，我们现正式开启 \textbf{天使轮} 融资。本轮融资将主要用于完成软硬件系统的高度产品化、底层SE安全芯片的联合流片以及首批政企标杆客户的样板工程落地。

\begin{table}[H]
	\centering
	\caption{天使轮融资计划}
	\label{tab:financing_plan}
	\renewcommand{\arraystretch}{1.4}
	\begin{tabular}{p{4cm}p{10cm}}
		\toprule
		\textbf{融资要素} & \textbf{详细说明} \\
		\midrule
		\textbf{融资轮次} & 天使轮 (Angel Round) \\
		\textbf{目标融资金额} & \textbf{1,000万 人民币} \\
		\textbf{出让股权比例} & \textbf{10\% - 15\%} \\
		\textbf{投前估值} & 约 6,000万 - 9,000万 人民币 \\
		\textbf{当前阶段} & 已完成算法底层重构与实验室Demo，初步达成芯片厂商合作意向 \\
		\bottomrule
	\end{tabular}
\end{table}

\begin{table}[H]
	\centering
	\caption{资金用途预算规划}
	\label{tab:fund_usage}
	\renewcommand{\arraystretch}{1.4}
	\begin{tabular}{p{3cm}cp{8.5cm}}
		\toprule
		\textbf{资金用途} & \textbf{占比} & \textbf{具体使用说明} \\
		\midrule
		\textbf{研发与芯片流片} & \textbf{40\%} & 支付流片代工（NRE）及掩膜费用，迭代嵌入式算法底层及硬件协处理器加速模块。 \\
		\textbf{软件平台开发} & \textbf{25\%} & 研发量子安全密钥管理平台（KMS）、敏捷性引擎架构及云端安全态势监控UI/UX。 \\
		\textbf{资质与合规认证} & \textbf{15\%} & 申报并获取国家密码管理局《商用密码产品认证证书》，推进 EAL4+ 安全认证。 \\
		\textbf{市场与标杆工程} & \textbf{10\%} & 在重点军工/政企园区落地2-3个规模化样板工程，开拓全国安防集成商渠道。 \\
		\textbf{日常运营储备} & \textbf{10\%} & 公司日常场地租赁、核心团队薪酬发放及法务/财务合规支出，保障至少18个月现金流。 \\
		\bottomrule
	\end{tabular}
\end{table}

